\section{논의}
\label{sec:discussion}

%======================================================================
\subsection{단계 위험을 좌우하는 요인}
\label{sec:discussion:drivers}

% P1: 학습 목적이 추정 대상을 정한다
\rev{학습 목적은 모델이 무엇을 추정할지를 정하며, 상대위험 기반 설계와 비교한 경로에서는 이 선택이 차이의 대부분을 만들었다.
상대위험 목적만으로는 단계마다 공통적인 위험 수준이 식별되지 않으며(\ref{sec:method:rq2}절), 실제로 본 적용판의 단계별 평균 예측은 거의 평평했고 단계별 기저위험을 결합하면 이 문제가 크게 줄었다. 그러나 목적을 바꾼 이득은 재보정의 영향을 받지 않는 단계 안·경기 안 비교에서도 남았는데, 이는 다음 단계라는 짧은 구간에 맞춘 목적이 남은 경기 전체의 탈락 순서와 다른 정보에 무게를 두기 때문으로 해석된다.
반대로 KRAFTON 적용판과 Cox 목적은 원래 과제에서 앞섰으므로, 설계는 성능 순위보다 추정 대상이 의사결정의 시간 단위와 맞는지로 비교해야 한다.}

% P2: 입력과 구조
\rev{입력에서는 특징의 수보다 줄어드는 구역과 주변 팀에 대한 팀의 위치를 담았는지가 중요하였다.
그런 정보가 없는 활동 점수에 서식지 맥락을 더하면 차이의 절반 이상이 메워졌고, 같은 정보를 일부 담은 KRAFTON 입력과는 예측력이 같았지만 합치면 성능이 올라 두 입력은 서로 다른 측면을 담는 것으로 보인다.
\rev{모델 구조가 같은 경기 안의 판별을 의미 있게 높이지 못한 것은}, 팀 사이의 관계가 경쟁 밀도와 관계 요약 특징으로 이미 입력에 요약되어 있기 때문일 수 있다.
반면 팀 축은 함께 살아 있는 팀 전체를 보므로 경기 단위의 맥락을 더해 단계 안 비교의 성능을 높였으며, 이 차이는 통합 AUC만으로는 보이지 않는다.}

% P3: 생태학적 해석
서식지 관점은 \rev{요인의 연관이} 단계에 따라 달라지는 양상을 해석하는 틀이 된다.
구역이 넓은 초반에는 주변 팀의 수가 위험을 크게 가르지만, 구역이 좁아진 후반에는 모든 팀이 가까이 모여 그 차이가 덜 중요해지고 남은 구역 안의 위치와 남은 팀원 수가 더 중요해지는 것으로 해석된다.
\rev{경계에 가까울수록 위험이 높고 그 차이가 구역이 작아질수록 커지는 양상은 서식지 가장자리의 개체가 더 큰 위험에 놓이는 가장자리 효과와 비슷하다~\cite{Ries04}.
이는 관측된 연관에 대한 해석이지만, 네 변수가 같은 단계 안에서도 판별력의 약 71\%를 재현했다는 점에서 이 틀은 위험을 요약하는 데 쓸모가 있다.}

%======================================================================
\subsection{활용과 한계}
\label{sec:discussion:implications}

% P4: 활용과 대규모 시공간 위험 예측에 대한 시사점
\rev{단계별 평균 위험 수준이 맞고 단계 안 위험 구간별 보정 오차도 기존 설계보다 작은 예측은 단계 사이의 위험을 같은 기준점에서 비교하게 하므로 경기 후 복기와 중계 분석에 쓰일 수 있고, 새 규칙에서 검증을 거친 뒤에는 구역 규칙 변경 전후의 비교에도 쓰일 수 있으며, 실시간 중계에는 실시간 데이터 공급이 더 필요하다.
이동하는 다수의 에이전트를 다루는 대규모 시공간 위험 예측에서도, 입력·생존 판정·관계에 같은 절단을 적용해 재구성으로 누수를 확인하고, 설계 요인의 효과를 통합 AUC뿐 아니라 단계 안·경기 안 지표와 단계별 보정으로 나누어 보고하며, 구조를 바꾸기 전에 추정 대상과 입력을 통제한 분해로 차이의 출처부터 확인해야 한다.}

% P5: 한계
\rev{결론은 다음 범위 안에서 읽어야 한다.
첫째, 입력은 모든 팀의 위치를 아는 관전자 관점에서 만들었고 예측 시점에 공개되는 다음 구역은 쓰지 않았으므로, 각 팀이 당시 알던 정보로 하는 의사결정 지원과는 구분된다.
둘째, 자료는 한 구역 규칙 아래 약 10주간의 스팀 서버 경기를 경기 단위로 무작위 분할한 것이고, 솔로 코호트는 하나뿐이며 일부 코호트에는 AI 계정이 많다.
셋째, 두 기존 설계는 같은 자료에 맞춘 적용판이며, 특히 KRAFTON 설계는 16팀 이하의 대회 경기용이고 일부 입력은 근사하였다.
넷째, RQ1의 계수는 관측 자료의 조건부 연관이며, 팀이 위치를 스스로 고르고 탈락 원인(구역 피해와 교전)을 구분하지 않았으므로 인과 효과로 읽을 수 없다.
다섯째, 구조에 관한 결론은 현재의 데이터 규모와 요약 특징, 코호트마다 하나의 초기화로 학습한 조건의 결과이다.}

%======================================================================
\section{결론}
\label{sec:conclusion}

% P6: 결론과 향후 과제
본 연구는 배틀로얄 경기의 대규모 텔레메트리에서 다음 단계 탈락 위험을 서식지 축소 관점의 단계 조건부 이산시간 hazard로 정식화하고, \rev{예측 시점 기준} 팀--단계 표본과 통제된 분해로 이 위험의 예측을 무엇이 좌우하는지 분석하였다.
\rev{경쟁 밀도, 이동 부담, 남은 팀원 수의 연관은 코호트에 걸쳐 일관되었고, \rev{Sagong 적용판과의 차이는 주로 서식지 맥락에서, KRAFTON 적용판과의 차이는 주로 정식화에서 왔으며,} 모델 구조는 같은 경기 안의 판별보다 경기 간 위험 차이에 작게 기여하였다.}
향후에는 각 팀이 관측할 수 있는 정보만 쓰는 의사결정 관점으로 모델을 확장하고, 수집 이후 적용된 구역 규칙 개편~\cite{PUBGBlueZone2026}을 자연 실험으로 삼아 위험 요인이 규칙이 바뀌어도 유지되는지 검증할 계획이다.

\section*{사사}
\rev{Anthropic의 생성형 AI인 Claude(Opus 5.5)를 분석·그림 코드의 작성과 초록 및 I--VI절의 교정에 사용하였으며, 저자는 AI가 만든 모든 코드, 수치, 문장을 검토하고 수정하였고 그 내용에 책임을 진다.}
% (연구비 사사 약 2줄: 게재 확정 후 추가. 이중 맹검 심사본에는 넣지 않는다.)
